\documentclass[10pt,a4paper]{amsart}
\usepackage[margin=19mm]{geometry}
\usepackage{amsmath,amssymb,mathtools}
\usepackage[scr=rsfs,cal=euler]{mathalfa}
\usepackage{xfrac}
\usepackage[unicode,hidelinks,bookmarks=false]{hyperref}
\newcommand{\ie}{i.e.\ }
\newcommand{\cf}{cf.\ }
\newcommand{\resp}{respectively\ }
\newcommand{\Subsection}{\subsection}
\providecommand{\nobreakdash}{\nobreak\hbox{-}}
\setlength{\emergencystretch}{2em}
\multlinegap0pt
\usepackage{amssymb}
\usepackage[shortlabels]{enumitem}
\newtheorem{theorem}{Theorem}
\newcommand{\Ham}{\operatorname{Ham}}
\newcommand{\id}{\operatorname{id}}
\title{GEOMETRIC PROPERTIES AND FLUX OF LOCALLY CONFORMALLY SYMPLECTIC DIFFEOMORPHISMS}
\author{S. Tchuiaga, F. Balibuno}
\date{Source: arXiv:2602.02368v1 (CC BY 4.0)}
\begin{document}
\maketitle

\noindent\emph{Source and licence.} GEOMETRIC PROPERTIES AND FLUX OF LOCALLY CONFORMALLY SYMPLECTIC DIFFEOMORPHISMS. Version arXiv:2602.02368v1 (CC BY 4.0), distributed by arXiv under the Creative Commons Attribution 4.0 International licence, http://creativecommons.org/licenses/by/4.0/, as recorded in the arXiv licence metadata for this exact version and shown on the abstract page https://arxiv.org/abs/2602.02368v1. Full licence notice: \texttt{licenses/arxiv-2602.02368-CC-BY-4.0.txt}.\par\smallskip

\noindent\textit{Editorial scope.} Counted unit: the article's theorem thm:deformation\_rigidity (source0.tex lines 592--594). The article's complete original proof of it is delivered with it (source0.tex lines 596--630, one \texttt{proof} environment of 35 lines). No mathematics is written, repaired or re-derived: the statement and the proof are the article's own lines, in the article's own notation, and no retained line is substituted or re-flowed.  No further source-local context is needed: every cross-reference of every retained line resolves inside the two delivered spans.  Nothing was omitted from any retained span, so the package carries no omission record.  No retained line contains a citation, so the wrapper carries no bibliography and no bibliography entry.  No retained line contains a figure environment, an image-inclusion command or drawing code, so no image is delivered and the package declares no asset dependency.  Editorial formatting only, in addition: an AMS \texttt{amsart} preamble that restates the article's own declarations and macros used by the retained lines, namely the environments \texttt{theorem} on separate counters, together with the standard packages those lines need.  The retained environments print in this document as Theorem 1 in order of appearance, whereas the article prints the same items with the numbering of its own full document.  The complete native source of the article as distributed in the arXiv e-print for this exact version is bundled unchanged as \texttt{sources/193848/source0.tex}.
	\begin{theorem}[Deformation Rigidity for Zero-Flux LCS Isotopies] \label{thm:deformation_rigidity}
		Let $(M, \Omega, \omega)$ be a compact LCS manifold and suppose that the LCS flux homomorphism $\widetilde{\Psi} : \widetilde{(\ker \Phi)}_0 \to H_\omega^1(M)$ has discrete period group $\Delta$ (Definition 3.3). Then the set of smooth isotopies $\{g_t\}_{t \in [0,1]} \subset (\ker \Phi)_0$ starting at the identity whose lift has vanishing flux, $\widetilde{\Psi}([\{g_t\}]) = 0$, is an open subset within the space of all smooth isotopies in $(\ker \Phi)_0$ starting at the identity (equipped with the $C^1$-topology on the path). Consequently, the Hamiltonian subgroup $\Ham_\Omega(M)$ is an open subgroup of $(\ker \Phi)_0$.
	\end{theorem}

	\begin{proof}
		We denote by
		\[
		\mathcal{P} = \{ \{g_t\}_{t \in [0,1]} \mid g_0 = \id_M \}
		\]
		the space of smooth isotopies in $(\ker \Phi)_0$, endowed with the $C^1$-topology.
		
		\textbf{Step 1: Construction of the Flux Map.} Define the map
		\[
		\alpha : \mathcal{P} \to \Omega^1(M), \quad \alpha(\{g_t\}) = \int_0^1 i_{X_t}\Omega \, dt,
		\]
		where $X_t$ is the time-dependent vector field generating the isotopy $\{g_t\}$. Since integration is continuous and $X_t$ depends continuously on $\{g_t\}$ in the $C^1$-topology, the map $\alpha$ is continuous. Next, let $\pi : \Omega^1(M) \to H_\omega^1(M)$, denote the natural projection. Then the composite map $\Psi_{\mathcal{P}} = \pi \circ \alpha : \mathcal{P} \to H_\omega^1(M)$, is continuous.
		
		\textbf{Step 2: The Zero-Flux Subset.} Define the subset of zero-flux isotopies as
		\[
		\mathcal{Z} = \{ \{g_t\} \in \mathcal{P} \mid \widetilde{\Psi}([\{g_t\}]) = 0 \}.
		\]
		In other words, an isotopy $\{g_t\}$ belongs to $\mathcal{Z}$ if
		\[
		\left[ \int_0^1 i_{X_t}\Omega \, dt \right] = 0 \quad \text{in } H_\omega^1(M).
		\]
		Let $\{g_t\}_0 \in \mathcal{Z}$ be an arbitrary zero-flux isotopy so that $\Psi_{\mathcal{P}}(\{g_t\}_0) = 0$.
		
		\textbf{Step 3: Local Stability via Continuity and Discreteness.} By the continuity of $\Psi_{\mathcal{P}}$, there is a $C^1$-neighborhood $U \subset \mathcal{P}$ of $\{g_t\}_0$ such that for every isotopy $\{\tilde{g}_t\} \in U$, the flux satisfies $\Psi_{\mathcal{P}}(\{\tilde{g}_t\}) \in V$, where $V$ is an open neighborhood of $0$ in $H_\omega^1(M)$. Since the lifted flux homomorphism $\widetilde{\Psi} : \widetilde{(\ker \Phi)}_0 \to H_\omega^1(M)$, has discrete period group $\Delta$, we can choose $V$ small enough so that
		\[
		V \cap \Delta = \{0\}.
		\]
		Because any isotopy (or its homotopy class) with fixed endpoints has its flux contained in $\Delta$, this forces the flux of every isotopy in $U$ to be zero. Thus, $\mathcal{Z}$ is open in $\mathcal{P}$.
		
		\textbf{Step 4: Openness of the Hamiltonian Subgroup.} Consider the endpoint evaluation map
		\[
		ev_1 : \mathcal{P} \to (\ker \Phi)_0, \quad ev_1(\{g_t\}) = g_1.
		\]
		This map is continuous and open in the $C^1$-topology. Since $\mathcal{Z}$ is an open subset of $\mathcal{P}$, its image $ev_1(\mathcal{Z})$ is an open subset of $(\ker \Phi)_0$. By definition, the Hamiltonian subgroup $\Ham_\Omega(M)$ consists of those diffeomorphisms that are endpoints of zero-flux isotopies, that is, $ev_1(\mathcal{Z}) \subseteq \Ham_\Omega(M)$. In particular, there exists an open neighborhood of the identity in $\Ham_\Omega(M)$, and since $\Ham_\Omega(M)$ is a subgroup, it follows that $\Ham_\Omega(M)$ is open in $(\ker \Phi)_0$.
	\end{proof}

\end{document}
